\section{Evaluación y laboratorio del agente LLM}

\subsection{Métricas cuantitativas de evaluación}

Antes de desplegar un agente, la empresa necesita construir un conjunto de métricas cuantitativas:

\begin{itemize}
    \item \textbf{Task Success Rate}: tasa de cumplimiento del objetivo (por ejemplo, organizar automáticamente un informe y obtener su aprobación)
    \item \textbf{Force Stop Rate}: proporción de conversaciones que el usuario interrumpe de forma activa
    \item \textbf{Tool Error Rate}: errores en la API o en las herramientas que provocan un workflow failure
    \item \textbf{Response Divergence}: la consistencia de las salidas entre distintas versiones del modelo
\end{itemize}

Estas métricas suelen organizarse por niveles de SLA: las Critical Task usan 99.9\% success, mientras que las Exploratory Task permiten un umbral más bajo.

\subsection{Laboratorios Human-in-the-loop}

Burak menciona que las empresas establecen un «Agent Lab»: testing con un grupo reducido de usuarios frente a escenarios reales. El proceso es el siguiente:

\begin{enumerate}
    \item Se simulan tareas reales y se deja que el agente se ejecute en un sistema tan realista como sea posible
    \item Se asigna un reviewer humano para hacer audit output de manera periódica y etiquetarlo
    \item Los resultados del label se convierten en un dashboard que dispara el retrain o el prompt tune
\end{enumerate}

\begin{knowledgebox}{El valor organizacional del Agent Lab}
El Agent Lab es el bridge between research and product: el equipo de investigación puede ver failure reales, el equipo de producto puede probar la nueva versión, y el equipo de gobernanza puede verificar si los guardrails de seguridad son eficaces.
\end{knowledgebox}

\subsection{Resumen de la sección}

El sistema de evaluación se compone de métricas cuantitativas y de laboratorios humanos, lo que garantiza que el agente sea eficiente y que, en caso de falla, se pueda localizar rápidamente el root cause.

